% questions.tex —— 2012 年·湖北卷·理科 · 题目内容片段（无答案）
% 由 tools/generate_projects.py 依来源总账自 2012/2012_数学真题_湖北理_试题.md 生成，勿手改（改源 MD 后重跑）


\begin{gknotice}
\textbf{来源说明}：本题面依据固定源 PDF 的两页 A3 题面逐题转录，原卷印刷页码为 737—738。第 737 面为选择题 1—10、填空题 11—16 及解答题 17；第 738 面为解答题 18—22。

\textbf{图形说明}：第 3 题二次函数图象、第 4 题三视图、第 8 题扇形与两个半圆、第 12 题程序框图、第 14 题双曲线与菱形、第 15 题圆与垂线、第 19 题图 1 与图 2，均由项目内 TikZ 矢量源重绘；其余各题原页无图。

\textbf{解析归属}：固定源不附答案或解析；配套答案与解析均由本库按本题面独立复算。
\end{gknotice}

\gkpart{一、选择题}

\begin{question}
\gkti{1}方程 $x^2+6x+13=0$ 的一个根是\mbox{（　　）}

（A）$-3+2\mathrm i$；（B）$3+2\mathrm i$；（C）$-2+3\mathrm i$；（D）$2+3\mathrm i$。

\end{question}
\begin{question}
\gkti{2}命题“$\exists x_0\in\complement_{\mathbf R}\mathbf Q,\ x_0^3\in\mathbf Q$”的否定是\mbox{（　　）}

（A）$\exists x_0\notin\complement_{\mathbf R}\mathbf Q,\ x_0^3\in\mathbf Q$；（B）$\exists x_0\in\complement_{\mathbf R}\mathbf Q,\ x_0^3\notin\mathbf Q$；（C）$\forall x\notin\complement_{\mathbf R}\mathbf Q,\ x^3\in\mathbf Q$；（D）$\forall x\in\complement_{\mathbf R}\mathbf Q,\ x^3\notin\mathbf Q$。

\Needspace{0.4\textheight}
\end{question}
\begin{question}
\gkti{3}已知二次函数 $y=f(x)$ 的图象如图所示，则它与 $x$ 轴所围成的图形的面积为\mbox{（　　）}
\nopagebreak[4]
\begin{center}% ============================================================
% f2012_hubei_science_q03_parabola.tex
% 2012 湖北卷（理科）第 3 题 二次函数 y=f(x) 的图象
% 图中抛物线开口向下，顶点 (0,1)，与 x 轴交于 (-1,0)、(1,0)
% 来源 PDF：普通高考/2012/2012湖北理.pdf（SHA-256 b84a88bc…07cc）
% 原页：第 1 页（A3 横排）第 1 栏，印面 737，第 3 题题干下方
% ============================================================
\begin{tikzpicture}[>=Stealth, line width=0.7pt, font=\small, x=1.4cm, y=1.4cm]
  \draw[->] (-2.3,0) -- (2.3,0) node[right] {$x$};
  \draw[->] (0,-1.55) -- (0,1.95) node[above] {$y$};
  \draw[thick, domain=-1.62:1.62, samples=120] plot (\x,{1-(\x)*(\x)});
  \node[above left] at (0,1) {$1$};
  \node at (-1.20,-0.20) {$-1$};
  \node at (1.20,-0.20) {$1$};
  \node[below left] at (0,0) {$O$};
\end{tikzpicture}
\end{center}
\nopagebreak[4]
（A）$\dfrac{2\pi}{5}$；（B）$\dfrac43$；（C）$\dfrac32$；（D）$\dfrac{\pi}{2}$。

\Needspace{0.4\textheight}
\end{question}
\begin{question}
\gkti{4}已知某几何体的三视图如图所示，则该几何体的体积为\mbox{（　　）}
\nopagebreak[4]
\begin{center}% ============================================================
% f2012_hubei_science_q04_views.tex
% 2012 湖北卷（理科）第 4 题 几何体三视图
% 正视图：直角梯形，底宽 2，左高 2，右高 4；侧视图：2×4 矩形内切半径 1 的圆；
% 俯视图：半径 1 的圆。几何体为半径 1、高 4 的圆柱被过轴心高度的斜平面所截。
% 来源 PDF：普通高考/2012/2012湖北理.pdf（SHA-256 b84a88bc…07cc）
% 原页：第 1 页（A3 横排）第 1 栏，印面 737，第 4 题题干下方
% ============================================================
\begin{tikzpicture}[>=Stealth, line width=0.7pt, font=\small, x=1cm, y=1cm]
  % ---------- 正视图 ----------
  \draw (0,0) -- (2,0) -- (2,4) -- (0,2) -- cycle;
  % 底宽 2
  \draw[<->] (0,-0.45) -- (2,-0.45);
  \draw (0,-0.58) -- (0,-0.32);
  \draw (2,-0.58) -- (2,-0.32);
  \node at (1,-0.85) {$2$};
  % 左高 2
  \draw[<->] (-0.45,0) -- (-0.45,2);
  \draw (-0.58,0) -- (-0.32,0);
  \draw (-0.58,2) -- (-0.32,2);
  \node at (-0.85,1) {$2$};
  % 右高 4
  \draw[<->] (2.45,0) -- (2.45,4);
  \draw (2.32,0) -- (2.58,0);
  \draw (2.32,4) -- (2.58,4);
  \node at (2.85,2) {$4$};
  \node at (1,-1.35) {正视图};

  % ---------- 侧视图 ----------
  \begin{scope}[xshift=4.1cm]
    \draw (0,0) rectangle (2,4);
    \draw (1,3) circle (1);
    \draw[<->] (0,-0.45) -- (2,-0.45);
    \draw (0,-0.58) -- (0,-0.32);
    \draw (2,-0.58) -- (2,-0.32);
    \node at (1,-0.85) {$2$};
    \node at (1,-1.35) {侧视图};
  \end{scope}

  % ---------- 俯视图 ----------
  \begin{scope}[yshift=-3.55cm]
    \draw (1,0) circle (1);
    \node at (1,-1.30) {俯视图};
  \end{scope}
\end{tikzpicture}
\end{center}
\nopagebreak[4]
（A）$\dfrac{8\pi}{3}$；（B）$3\pi$；（C）$\dfrac{10\pi}{3}$；（D）$6\pi$。

\end{question}
\begin{question}
\gkti{5}设 $a\in\mathbf Z$，且 $0\leqslant a<13$，若 $5^{2012}+a$ 能被 $13$ 整除，则 $a=$\mbox{（　　）}

（A）$0$；（B）$1$；（C）$11$；（D）$12$。

\end{question}
\begin{question}
\gkti{6}设 $a$，$b$，$c$，$x$，$y$，$z$ 是正数，且 $a^2+b^2+c^2=10$，$x^2+y^2+z^2=40$，$ax+by+cz=20$，则 $\dfrac{a+b+c}{x+y+z}=$\mbox{（　　）}

（A）$\dfrac14$；（B）$\dfrac13$；（C）$\dfrac12$；（D）$\dfrac34$。

\end{question}
\begin{question}
\gkti{7}定义在 $(-\infty,0)\cup(0,+\infty)$ 上的函数 $f(x)$，如果对于任意给定的等比数列 $\{a_n\}$，$\{f(a_n)\}$ 仍是等比数列，则称 $f(x)$ 为“保等比数列函数”。现有定义在 $(-\infty,0)\cup(0,+\infty)$ 上的如下函数：

① $f(x)=x^2$；② $f(x)=2^x$；③ $f(x)=\sqrt{|x|}$；④ $f(x)=\ln|x|$。

则其中是“保等比数列函数”的 $f(x)$ 的序号为\mbox{（　　）}

（A）①②；（B）③④；（C）①③；（D）②④。

\Needspace{0.4\textheight}
\end{question}
\begin{question}
\gkti{8}如图，在圆心角为直角的扇形 $OAB$ 中，分别以 $OA$，$OB$ 为直径作两个半圆。在扇形 $OAB$ 内随机取一点，则此点取自阴影部分的概率是\mbox{（　　）}
\nopagebreak[4]
\begin{center}% ============================================================
% f2012_hubei_science_q08_sector.tex
% 2012 湖北卷（理科）第 8 题 直角扇形 OAB 内的两个半圆
% 扇形圆心角 90°，OA 水平、OB 竖直；以 OA、OB 为直径各作一个半圆。
% 阴影＝两半圆的公共部分（透镜形）＋扇形内两半圆之外的部分（月牙形）。
% 来源 PDF：普通高考/2012/2012湖北理.pdf（SHA-256 b84a88bc…07cc）
% 原页：第 1 页（A3 横排）第 2 栏上部，印面 737，第 8 题题干下方
% ============================================================
\begin{tikzpicture}[>=Stealth, line width=0.7pt, font=\small, scale=2.3]
  \def\r{1}
  \coordinate (O) at (0,0);
  \coordinate (A) at (\r,0);
  \coordinate (B) at (0,\r);

  % 扇形内先铺阴影，再用白色挖去两个半圆，最后补回公共部分
  \begin{scope}
    \clip (O) -- (A) arc (0:90:\r) -- cycle;
    \fill[black!15] (O) -- (A) arc (0:90:\r) -- cycle;
    \fill[white] (\r/2,0) circle (\r/2);
    \fill[white] (0,\r/2) circle (\r/2);
  \end{scope}
  \begin{scope}
    \clip (\r/2,0) circle (\r/2);
    \clip (0,\r/2) circle (\r/2);
    \fill[black!15] (-0.2,-0.2) rectangle (\r+0.2,\r+0.2);
  \end{scope}

  % 轮廓
  \draw (O) -- (A);
  \draw (O) -- (B);
  \draw (A) arc (0:90:\r);
  \draw (O) arc (180:0:\r/2);
  \draw (B) arc (90:-90:\r/2);

  \node[below left] at (O) {$O$};
  \node[below] at (A) {$A$};
  \node[left] at (B) {$B$};
\end{tikzpicture}
\end{center}
\nopagebreak[4]
（A）$1-\dfrac{2}{\pi}$；（B）$\dfrac12-\dfrac{1}{\pi}$；（C）$\dfrac{2}{\pi}$；（D）$\dfrac{1}{\pi}$。

\end{question}
\begin{question}
\gkti{9}函数 $f(x)=x\cos x^2$ 在区间 $[0,4]$ 上的零点个数为\mbox{（　　）}

（A）$4$；（B）$5$；（C）$6$；（D）$7$。

\end{question}
\begin{question}
\gkti{10}我国古代数学名著《九章算术》中“开立圆术”曰：置积尺数，以十六乘之，九而一，所得开立方除之，即立圆径。“开立圆术”相当于给出了已知球的体积 $V$，求其直径 $d$ 的一个近似公式 $d\approx\sqrt[3]{\dfrac{16}{9}V}$。人们还用过一些类似的近似公式。根据 $\pi=3.14159\cdots$ 判断，下列近似公式中最精确的一个是\mbox{（　　）}

（A）$d\approx\sqrt[3]{\dfrac{16}{9}V}$；（B）$d\approx\sqrt[3]{2V}$；（C）$d\approx\sqrt[3]{\dfrac{300}{157}V}$；（D）$d\approx\sqrt[3]{\dfrac{21}{11}V}$。

\end{question}
\gkpart{二、填空题}

\begin{question}
\gkti{11}设 $\triangle ABC$ 的内角 $A$，$B$，$C$ 所对的边分别为 $a$，$b$，$c$。若 $(a+b-c)(a+b+c)=ab$，则角 $C=$\gkblank{}。

\Needspace{0.4\textheight}
\end{question}
\begin{question}
\gkti{12}阅读如图所示的程序框图，运行相应的程序，输出的结果 $s=$\gkblank{}。
\nopagebreak[4]
\begin{center}% ============================================================
% f2012_hubei_science_q12_flowchart.tex
% 2012 湖北卷（理科）第 12 题 程序框图
% 初值 a=1、s=0、n=1；循环体 s=s+a、a=a+2；判断 n<3，是则 n=n+1 返回，
% 否则输出 s。运行结果 s=9。
% 来源 PDF：普通高考/2012/2012湖北理.pdf（SHA-256 b84a88bc…07cc）
% 原页：第 1 页（A3 横排）第 3 栏，印面 737，第 12 题题干下方
% ============================================================
\begin{tikzpicture}[>=Stealth, line width=0.7pt, font=\small]
  % 开始
  \draw[rounded corners=8pt] (-0.55,0.50) rectangle (0.55,0.10);
  \node at (0,0.30) {开始};
  \draw[->] (0,0.10) -- (0,-0.30);

  % 初值
  \draw (-1.55,-0.30) rectangle (1.55,-0.80);
  \node at (0,-0.55) {$a=1,\ s=0,\ n=1$};
  \draw[->] (0,-0.80) -- (0,-1.20);

  % s = s + a
  \draw (-0.85,-1.20) rectangle (0.85,-1.70);
  \node at (0,-1.45) {$s=s+a$};
  \draw[->] (0,-1.70) -- (0,-2.05);

  % a = a + 2
  \draw (-0.85,-2.05) rectangle (0.85,-2.55);
  \node at (0,-2.30) {$a=a+2$};
  \draw[->] (0,-2.55) -- (0,-2.85);

  % 判断 n < 3
  \draw (0,-2.85) -- (1.05,-3.55) -- (0,-4.25) -- (-1.05,-3.55) -- cycle;
  \node at (0,-3.55) {$n<3$?};

  % 是：n = n + 1，回边
  \draw[->] (0,-4.25) -- (0,-4.60);
  \node[right] at (0.08,-4.38) {是};
  \draw (-0.80,-4.60) rectangle (0.80,-5.10);
  \node at (0,-4.85) {$n=n+1$};
  \draw (-0.80,-4.85) -- (-2.15,-4.85) -- (-2.15,-1.45);
  \draw[->] (-2.15,-1.45) -- (-0.85,-1.45);

  % 否：输出 s
  \draw[->] (1.05,-3.55) -- (2.15,-3.55);
  \node[above] at (1.30,-3.50) {否};
  \draw (2.45,-3.30) -- (3.95,-3.30) -- (3.45,-3.80) -- (1.95,-3.80) -- cycle;
  \node at (2.95,-3.55) {输出 $s$};
  \draw[->] (2.95,-3.80) -- (2.95,-4.50);

  % 结束
  \draw[rounded corners=8pt] (2.35,-4.50) rectangle (3.55,-4.95);
  \node at (2.95,-4.72) {结束};
\end{tikzpicture}
\end{center}
\nopagebreak[4]
\end{question}
\begin{question}
\gkti{13}回文数是指从左到右读与从右到左读都一样的正整数。如 $22$，$121$，$3443$，$94249$ 等。显然 $2$ 位回文数有 $9$ 个：$11$，$22$，$33$，$\cdots$，$99$；$3$ 位回文数有 $90$ 个：$101$，$111$，$121$，$\cdots$，$191$，$202$，$\cdots$，$999$。则

（1） 4 位回文数有\gkblank{}个；

（2） $2n+1$（$n\in\mathbf N^*$）位回文数有\gkblank{}个。

\Needspace{0.4\textheight}
\end{question}
\begin{question}
\gkti{14}如图，双曲线 $\dfrac{x^2}{a^2}-\dfrac{y^2}{b^2}=1$（$a,b>0$）的两顶点为 $A_1$，$A_2$，虚轴两端点为 $B_1$，$B_2$，两焦点为 $F_1$，$F_2$。若以 $A_1A_2$ 为直径的圆内切于菱形 $F_1B_1F_2B_2$，切点分别为 $A$，$B$，$C$，$D$。则

（1） 双曲线的离心率 $e=$\gkblank{}；

（2） 菱形 $F_1B_1F_2B_2$ 的面积 $S_1$ 与矩形 $ABCD$ 的面积 $S_2$ 的比值 $\dfrac{S_1}{S_2}=$\gkblank{}。
\nopagebreak[4]
\begin{center}% ============================================================
% f2012_hubei_science_q14_hyperbola.tex
% 2012 湖北卷（理科）第 14 题 双曲线与内切圆、菱形、矩形
% 双曲线 x^2/a^2-y^2/b^2=1，顶点 A1、A2，虚轴端点 B1、B2，焦点 F1、F2；
% 以 A1A2 为直径的圆（半径 a）内切于菱形 F1B1F2B2，切点 A、B、C、D 成矩形。
% 图中取 a=1.6，e=(1+sqrt5)/2，b=a*sqrt(e^2-1)=2.035，c=ae=2.589。
% 来源 PDF：普通高考/2012/2012湖北理.pdf（SHA-256 b84a88bc…07cc）
% 原页：第 1 页（A3 横排）第 3 栏上部，印面 737，第 14 题题干下方
% ============================================================
\begin{tikzpicture}[>=Stealth, line width=0.7pt, font=\small, x=1cm, y=1cm]
  \def\a{1.6}
  \def\b{2.035}
  \def\c{2.589}
  \def\xt{0.989}
  \def\yt{1.258}

  \draw[->] (-3.55,0) -- (3.75,0) node[right] {$x$};
  \draw[->] (0,-3.35) -- (0,3.45) node[above] {$y$};

  % 双曲线两支
  \draw[thick, domain=-2.55:2.55, samples=120] plot ({1.6*sqrt(1+(\x*\x)/4.14103)},\x);
  \draw[thick, domain=-2.55:2.55, samples=120] plot ({-1.6*sqrt(1+(\x*\x)/4.14103)},\x);

  % 菱形
  \draw (-\c,0) -- (0,\b) -- (\c,0) -- (0,-\b) -- cycle;

  % 内切圆
  \draw (\a,0) arc (0:360:\a);

  % 矩形 ABCD
  \draw (-\xt,\yt) -- (\xt,\yt) -- (\xt,-\yt) -- (-\xt,-\yt) -- cycle;

  % 切点
  \fill (\xt,\yt) circle (1.2pt) node[right] {$A$};
  \fill (-\xt,\yt) circle (1.2pt) node[left] {$B$};
  \fill (-\xt,-\yt) circle (1.2pt) node[left] {$C$};
  \fill (\xt,-\yt) circle (1.2pt) node[right] {$D$};

  % 顶点与焦点
  \fill (-\a,0) circle (1.2pt) node[below left] {$A_1$};
  \fill (\a,0) circle (1.2pt) node[below right] {$A_2$};
  \fill (0,\b) circle (1.2pt) node[right] {$B_2$};
  \fill (0,-\b) circle (1.2pt) node[right] {$B_1$};
  \fill (-\c,0) circle (1.2pt) node[below left] {$F_1$};
  \fill (\c,0) circle (1.2pt) node[below right] {$F_2$};
  \node[below left] at (0,0) {$O$};
\end{tikzpicture}
\end{center}
\nopagebreak[4]
\Needspace{0.4\textheight}
\end{question}
\begin{question}
\gkti{15}如图，点 $D$ 在 $\odot O$ 的弦 $AB$ 上移动，$AB=4$，连接 $OD$，过点 $D$ 作 $OD$ 的垂线交 $\odot O$ 于点 $C$，则 $CD$ 的最大值为\gkblank{}。
\nopagebreak[4]
\begin{center}% ============================================================
% f2012_hubei_science_q15_circle.tex
% 2012 湖北卷（理科）第 15 题 圆内定长弦与垂线
% 圆 O 半径 2.4，弦 AB 长 4（A、B 在圆上，弦垂直于 x 轴方向摆放）；
% D 在弦 AB 上，连 OD，过 D 作 OD 的垂线交圆于 C。
% 来源 PDF：普通高考/2012/2012湖北理.pdf（SHA-256 b84a88bc…07cc）
% 原页：第 1 页（A3 横排）第 3 栏中部，印面 737，第 15 题题干下方
% ============================================================
\begin{tikzpicture}[>=Stealth, line width=0.7pt, font=\small, x=1cm, y=1cm]
  \coordinate (O) at (0,0);
  \coordinate (A) at (-1.327,-2);
  \coordinate (B) at (-1.327,2);
  \coordinate (D) at (-1.327,1.6);
  \coordinate (C) at (-0.403,2.366);

  \draw (0,0) circle (2.4);
  \draw (A) -- (B);
  \draw (O) -- (D);
  \draw (D) -- (C);

  % 直角标记（D 处，OD 与 DC 之间）
  \draw (-1.212,1.461) -- (-1.073,1.576) -- (-1.188,1.715);

  \fill (O) circle (1.2pt) node[below] {$O$};
  \fill (A) circle (1.2pt) node[below left] {$A$};
  \fill (B) circle (1.2pt) node[above left] {$B$};
  \fill (D) circle (1.2pt) node[left] {$D$};
  \fill (C) circle (1.2pt);
  \node at (-0.06,2.62) {$C$};
\end{tikzpicture}
\end{center}
\nopagebreak[4]
\end{question}
\begin{question}
\gkti{16}在直角坐标系 $xOy$ 中，以原点 $O$ 为极点，$x$ 轴的正半轴为极轴建立极坐标系。已知射线 $\theta=\dfrac{\pi}{4}$ 与曲线 $\begin{cases}x=t+1,\\y=(t-1)^2\end{cases}$（$t$ 为参数）相交于 $A$，$B$ 两点，则线段 $AB$ 的中点的直角坐标为\gkblank{}。

\end{question}
\gkpart{三、解答题}

\begin{question}
\gkti{17}已知向量 $\boldsymbol a=(\cos\omega x-\sin\omega x,\ \sin\omega x)$，$\boldsymbol b=(-\cos\omega x-\sin\omega x,\ 2\sqrt3\cos\omega x)$，设函数 $f(x)=\boldsymbol a\cdot\boldsymbol b+\lambda$（$x\in\mathbf R$）的图象关于直线 $x=\pi$ 对称，其中 $\omega$，$\lambda$ 为常数，且 $\omega\in\left(\dfrac12,1\right)$。

（1） 求函数 $f(x)$ 的最小正周期；

（2） 若 $y=f(x)$ 的图象经过点 $\left(\dfrac{\pi}{4},0\right)$，求函数 $f(x)$ 在区间 $\left[0,\dfrac{3\pi}{5}\right]$ 上的取值范围。

\end{question}
\begin{question}
\gkti{18}已知等差数列 $\{a_n\}$ 前三项的和为 $-3$，前三项的积为 $8$。

（1） 求等差数列 $\{a_n\}$ 的通项公式；

（2） 若 $a_2$，$a_3$，$a_1$ 成等比数列，求数列 $\{|a_n|\}$ 的前 $n$ 项和。

\Needspace{0.4\textheight}
\end{question}
\begin{question}
\gkti{19}如图 1，$\angle ACB=45^\circ$，$BC=3$，过动点 $A$ 作 $AD\perp BC$，垂足 $D$ 在线段 $BC$ 上且异于点 $B$，连接 $AB$，沿 $AD$ 将 $\triangle ABD$ 折起，使 $\angle BDC=90^\circ$（如图 2 所示）。

（1） 当 $BD$ 的长为多少时，三棱锥 $A-BCD$ 的体积最大；

（2） 当三棱锥 $A-BCD$ 的体积最大时，设点 $E$，$M$ 分别为棱 $BC$，$AC$ 的中点，试在棱 $CD$ 上确定一点 $N$，使得 $EN\perp BM$，并求 $EN$ 与平面 $BMN$ 所成角的大小。
\nopagebreak[4]
\begin{center}% ============================================================
% f2012_hubei_science_q19_fig1.tex
% 2012 湖北卷（理科）第 19 题 图 1（折起前的平面图）
% B、D、C 共线，AD 垂直 BC，∠ACB=45°，BC=3（图中 BD=1，DC=AD=2）。
% 来源 PDF：普通高考/2012/2012湖北理.pdf（SHA-256 b84a88bc…07cc）
% 原页：第 2 页（A3 横排）第 1 栏下部，印面 738，第 19 题题干下方
% ============================================================
\begin{tikzpicture}[>=Stealth, line width=0.7pt, font=\small, x=1cm, y=1cm]
  \coordinate (B) at (0,0);
  \coordinate (D) at (1,0);
  \coordinate (C) at (3,0);
  \coordinate (A) at (1,2);

  \draw (A) -- (B) -- (C) -- cycle;
  \draw (A) -- (D);

  \fill (A) circle (1.2pt) node[above] {$A$};
  \fill (B) circle (1.2pt) node[below left] {$B$};
  \fill (C) circle (1.2pt) node[below right] {$C$};
  \fill (D) circle (1.2pt) node[below] {$D$};
\end{tikzpicture}
\end{center}
\nopagebreak[4]
\nopagebreak[4]
\begin{center}% ============================================================
% f2012_hubei_science_q19_fig2.tex
% 2012 湖北卷（理科）第 19 题 图 2（沿 AD 折起后的三棱锥 A-BCD）
% AD ⊥ 底面 BCD，∠BDC=90°；M 为 AC 中点，E 为 BC 中点；N 待定，原页未画出。
% 采用固定斜投影：D=(1.00,1.05)、B=(0.10,0.90)、C=(2.60,0.95)、A=(1.00,3.75)，
% 其中 E、M 由相应顶点直接取中点得到。
% 来源 PDF：普通高考/2012/2012湖北理.pdf（SHA-256 b84a88bc…07cc）
% 原页：第 2 页（A3 横排）第 1 栏下部，印面 738，第 19 题题干下方
% ============================================================
\begin{tikzpicture}[>=Stealth, line width=0.7pt, font=\small, x=1cm, y=1cm]
  \coordinate (A) at (1.00,3.75);
  \coordinate (B) at (0.10,0.90);
  \coordinate (C) at (2.60,0.95);
  \coordinate (D) at (1.00,1.05);
  \coordinate (E) at (1.35,0.925);
  \coordinate (M) at (1.80,2.35);

  % 可见棱
  \draw (A) -- (B);
  \draw (A) -- (C);
  \draw (B) -- (C);
  \draw (B) -- (M);

  % 被遮棱
  \draw[dashed] (A) -- (D);
  \draw[dashed] (D) -- (B);
  \draw[dashed] (D) -- (C);

  % D 处直角标记（DA 与 DC 之间）
  \draw (1.00,1.25) -- (1.20,1.2375) -- (1.20,1.0375);

  \fill (A) circle (1.2pt) node[above] {$A$};
  \fill (B) circle (1.2pt) node[below left] {$B$};
  \fill (C) circle (1.2pt) node[right] {$C$};
  \fill (D) circle (1.2pt) node[below left] {$D$};
  \fill (E) circle (1.2pt) node[below right] {$E$};
  \fill (M) circle (1.2pt) node[right] {$M$};
\end{tikzpicture}
\end{center}
\nopagebreak[4]
\end{question}
\begin{question}
\gkti{20}根据以往的经验，某工程施工期间的降水量 $X$（单位：mm）对工期的影响如下表：

\begin{center}
\begin{tabular}{lllll}
\toprule
降水量 $X$ & $X<300$ & $300\leqslant X<700$ & $700\leqslant X<900$ & $X\geqslant900$ \\
\midrule
工期延误天数 $Y$ & $0$ & $2$ & $6$ & $10$ \\
\bottomrule
\end{tabular}
\end{center}

历年气象资料表明，该工程施工期间降水量 $X$ 小于 $300$，$700$，$900$ 的概率分别为 $0.3$，$0.7$，$0.9$。求：

（1） 工期延误天数 $Y$ 的均值与方差；

（2） 在降水量 $X$ 至少是 $300$ 的条件下，工期延误不超过 $6$ 天的概率。

\end{question}
\begin{question}
\gkti{21}设 $A$ 是单位圆 $x^2+y^2=1$ 上的任意一点，$l$ 是过点 $A$ 与 $x$ 轴垂直的直线，$D$ 是直线 $l$ 与 $x$ 轴的交点，点 $M$ 在直线 $l$ 上，且满足 $|DM|=m|DA|$（$m>0$ 且 $m\neq1$）。当点 $A$ 在圆上运动时，记点 $M$ 的轨迹为曲线 $C$。

（1） 求曲线 $C$ 的方程，判断曲线 $C$ 为何种圆锥曲线，并求其焦点坐标；

（2） 过原点且斜率为 $k$ 的直线交曲线 $C$ 于 $P$，$Q$ 两点，其中 $P$ 在第一象限，它在 $y$ 轴上的射影为点 $N$，直线 $QN$ 交曲线 $C$ 于另一点 $H$。是否存在 $m$，使得对任意的 $k>0$，都有 $PQ\perp PH$？若存在，求 $m$ 的值；若不存在，请说明理由。

\end{question}
\begin{question}
\gkti{22}（1） 已知函数 $f(x)=rx-x^r+(1-r)$（$x>0$），其中 $r$ 为有理数，且 $0<r<1$，求 $f(x)$ 的最小值；

（2） 试用（1）的结果证明如下命题：设 $a_1\geqslant0$，$a_2\geqslant0$，$b_1$，$b_2$ 为正有理数。若 $b_1+b_2=1$，则 $a_1^{b_1}a_2^{b_2}\leqslant a_1b_1+a_2b_2$；

（3） 请将（2）中的命题推广到一般形式，并用数学归纳法证明你所推广的命题。

注：当 $\alpha$ 为正有理数时，有求导公式 $(x^\alpha)'=\alpha x^{\alpha-1}$。

\end{question}